\chapter[Referencial teórico]{Referencial teórico}
\label{cap:referencial}

% {aqui vai um parágrafo de abertura explicando quais conceitos o capítulo apresenta e por quê, porque conecta o capítulo à questão de pesquisa}
%
% Notas do doc: Engenharia de Software Experimental; necessário (histórias de usuário e CA, LLMs);
% Métodos ágeis; Modelo de qualidade (ISO); Estatística.

\section{Engenharia de Software Experimental}
\label{sec:es-experimental}
% {aqui vão o conceito de experimentação em ES e os tipos de estudo empírico (experimento, estudo de caso, survey, revisão), sustenta a revisão e o estudo empírico; base: Wohlin et al.}

\subsection{Paradigma GQM}
\label{sec:ref-gqm}
% {aqui vai o Goal-Question-Metric e o template de objetivo, porque nosso objetivo e as QPs foram definidos com GQM adaptado}

\subsection{Revisão sistemática da literatura}
\label{sec:ref-rsl}
% {aqui vão o conceito de RSL, a estrutura PICOC e as diretrizes PRISMA 2020, porque fundamentam o método do capítulo da revisão}

\subsection{Experimento}
\label{sec:ref-experimento}
% {aqui vão as fases do experimento: definição do escopo, planejamento, operação, análise e interpretação, apresentação e empacotamento, porque fundamenta o capítulo de planejamento}

% [necessário]
\section{Histórias de usuário e critérios de aceitação}
\label{sec:hu-ca}
% {aqui vai a introdução da seção ligando histórias de usuário e critérios, porque é a população e o desfecho do PICOC}

\subsection{Histórias de usuário}
% {aqui vão a definição, o formato "Como... quero... para..." e o modelo INVEST, porque é a entrada da geração}

\subsection{Critérios de aceitação}
% {aqui vão a definição, o papel na aceitação da história, os formatos (lista de regras, cenários) e características de bons critérios (testável, atômico, não ambíguo), porque é o artefato central do trabalho}

\section{Modelos de linguagem de grande porte}
\label{sec:llm}
% {aqui vão o conceito de LLM (transformer, pré-treinamento), modelos proprietários e abertos e formas de hospedagem (nuvem pública, privada, local), porque é a intervenção do PICOC e base da QP4}

\section{Métodos ágeis e engenharia de requisitos ágil}
\label{sec:ageis}
% {aqui vão os princípios ágeis e como requisitos são tratados (backlog, refinamento, Definition of Done), porque o contexto do GQM são aplicações web ágeis; pode ser curta}

\section{\textit{Behavior-Driven Development} e \textit{Gherkin}}
\label{sec:bdd}
% {aqui vão o BDD, a sintaxe Gherkin (Feature, Scenario, Given/When/Then) com exemplo e ferramentas (Cucumber, Behave), porque Gherkin é o formato de saída no objetivo GQM}

\section{Técnicas de apoio à geração com LLMs}
\label{sec:tecnicas-llm}
% {aqui vão prompting (zero-shot, few-shot, chain-of-thought), RAG, fine-tuning, entrada multimodal e agentes, porque as QP5 e QP6 classificam os estudos por essas técnicas; definir antes evita explicar no meio dos resultados}
